\documentclass[11pt]{article}
\usepackage[margin=1in]{geometry}
\usepackage{amsmath,amssymb}
\usepackage{tikz}
\usepackage{float}
\usepackage[hidelinks]{hyperref}

\setlength{\parindent}{0pt}
\setlength{\parskip}{0.7\baselineskip}

\begin{document}

\section{CP and ILP via Graph Burning}

We will be working through some advanced work on CP and ILP modelling using the graph burning problem as a central theme. We will look at several papers on the topic, and compare different approaches looking at the development over time.

Here are links to the papers that we will be looking at, with a brief description of what we will discuss in each:

\subsection{A survey on graph burning in general}
\url{https://arxiv.org/abs/2009.10642}
\begin{itemize}
    \item here we will talk about the definition of the graph burning problem, and mention the graph burning conjecture
\end{itemize}

\subsection{CP and ILP formulations}
\url{https://www.mdpi.com/2227-7390/10/15/2777}
\begin{itemize}
    \item here we will look in fair detail at each of the CSP/ILP models for graph burning, and discuss how they work, how they differ and where they are similar.
    \item we will talk about the interaction with binary search
    \item we will look at the relative performances, and what the authors think about why these perform as they do
\end{itemize}

\subsection{Parameterised complexity (this in less detail)}
\url{https://arxiv.org/pdf/2007.08811}
\begin{itemize}
    \item we will talk briefly about what parameterised algorithms are, and why and when they help us
    \item we will talk very briefly about why we would care about these algorithms
\end{itemize}

\subsection{A greedy heuristic algorithm (less detail)}
\url{https://arxiv.org/abs/2401.07577}
\begin{itemize}
    \item we will talk about what a greedy algorithm is, what a heuristic algorithm is, and what they can and cannot do
    \item we will talk about the differences (and advantages and disadvantages) of heuristic, parameterised, and exact exponential algorithms
    \item we will very briefly mention the link to approximation algorithms, and what an approximation algorithm is
\end{itemize}

\subsection{An exact row-generation algorithm (integer programming)}
\url{https://arxiv.org/abs/2404.17080}
\begin{itemize}
    \item here we will look at the updated ILP approach, and compare it briefly to the ones in the earlier paper
    \item we will look at the results here and talk about the comparison to other approaches
\end{itemize}

\end{document}
